% Chapter 4: Formal State Model
% Contains: Device Identity, Network Topology, Cost Model, HLC, CRDTs

\section{Formal State Model}
% ============================================================

This section provides the formal definitions underlying Avena's distributed state management.

\subsection{Device Identity}

Let $\mathbf{D}$ denote the set of Avena devices. Each device $d \in \mathbf{D}$ possesses:
\begin{itemize}
    \item A keypair $(sk_d, pk_d)$ where $sk_d$ is an Ed25519 private key and $pk_d$ is the corresponding public key
    \item Identity $\text{id}(d) = H(pk_d)$ where $H$ is SHA-256 truncated to 128 bits (16 bytes), encoded as lowercase Base32 (RFC 4648, no padding)
    \item A certificate $\text{cert}(d)$ binding the public key to authorized capabilities
\end{itemize}

Devices form a containment hierarchy. Define the relation $d_1 \prec d_2$ (``$d_1$ is contained in $d_2$'') when $d_1$ is a nested Avena instance spawned by $d_2$. The transitive closure $\prec^*$ defines ancestral containment.

Each device $d$ has associated metadata:
\begin{itemize}
    \item $\text{roles}(d) \subseteq \Sigma_r$, a set of role labels
    \item $\text{labels}(d) : \Sigma_k \to \Sigma_v$, a partial function from label keys to values
    \item $\text{scope}(d) \subseteq \mathcal{P}(\Sigma_s)$, permitted message subjects as a set of patterns
\end{itemize}

\textbf{Capability attenuation:} For nested devices, $\text{scope}(d_1) \subseteq \text{scope}(d_2)$ whenever $d_1 \prec d_2$.

\subsection{Network Topology}

The network is modeled as a time-varying dual-channel graph $G(t) = (D, E_h(t), E_l(t))$ where:
\begin{itemize}
    \item $E_h(t) \subseteq D \times D$ is the set of active high-rate links (Wireguard tunnels) at time $t$
    \item $E_l(t) \subseteq D \times D$ is the set of active low-rate links (LoRa) at time $t$
\end{itemize}

\begin{definition}[Message Relay]
Device $d$ will relay application message $m$ for interest $i$ if:
\begin{itemize}
    \item $\text{relay}(d) = \text{None} \Rightarrow d$ is the interest owner
    \item $\text{relay}(d) = \text{Priority}(p) \Rightarrow \text{priority}(i) \geq p$
    \item $\text{relay}(d) = \text{Full} \Rightarrow$ always relay
    \item $\text{relay}(d) = \text{Subjects}(P) \Rightarrow \text{pattern}(i)$ matches some pattern in $P$
\end{itemize}
\end{definition}

\begin{definition}[$T$-Connectivity]
The network is $(T, i)$-connected for interest $i$ if for all device pairs (publisher, subscriber) where the subscriber holds $i$, and all times $\tau_0$, there exists a temporal admissible path completing within time $T$.
\end{definition}

\subsection{Cost Model}
\label{sec:costmodel}

Each device $d$ maintains cost metadata $c(d)$ propagated via gossip:

\begin{lstlisting}[language=Rust]
struct CostMetadata {
    device_id: DeviceId,
    timestamp: HLC,
    
    // Node state
    battery_fraction: Option<f32>,  // (0.0, 1.0] or None for wall-powered
    storage_available_mb: u32,
    queue_depth: u32,
    
    // Per-neighbor link costs
    links: HashMap<DeviceId, LinkCost>,
}

struct LinkCost {
    latency_ms: u32,
    bandwidth_kbps: u32,
    loss_percent: f32,
    dollar_cost_per_mb: Option<u32>,
    energy_cost_per_mb: Option<u32>,
}
\end{lstlisting}

Note: \texttt{battery\_fraction} is bounded to be strictly positive (the interval $(0.0, 1.0]$) because a device with zero battery cannot transmit messages.

\textbf{Cost function:} For message $m$ with interest $i$ traversing edge $(d_1, d_2)$:

\begin{equation}
\text{cost}(m, i, d_1, d_2) = \sum_j w_j \cdot f_j(m, i, d_1, d_2, c(d_1), c(d_2))
\end{equation}

Component functions:
\begin{align}
f_{\text{latency}} &= \text{latency}(d_1, d_2) \\
f_{\text{dollar}} &= \text{dollar\_cost}(d_1, d_2) \cdot \text{size}(m) \\
f_{\text{energy}} &= \frac{\text{energy\_cost}(d_1, d_2) \cdot \text{size}(m)}{\text{battery}(d_1)} \\
f_{\text{queue}} &= \frac{\text{queue\_depth}(d_1)}{\text{bandwidth}(d_1, d_2)} \\
f_{\text{priority}} &= (\text{max\_priority} - \text{priority}(i)) \cdot \text{latency}(d_1, d_2)
\end{align}

Weights $w_j$ are configurable at network level (fleet policy) and device level (local constraints). When cost metadata is stale or unavailable, routing falls back to Babel shortest path.

\subsubsection{Cost Profiles}

Rather than requiring manual weight tuning, Avena provides preset profiles for common deployment scenarios:

\begin{lstlisting}[language=Rust]
enum CostProfile {
    Balanced,          // Default: equal weight to all factors
    EnergyConstrained, // Solar/battery: prioritize energy conservation
    CostSensitive,     // Metered cellular: minimize dollar cost
    LatencyCritical,   // V2X/real-time: minimize latency, ignore cost
}
\end{lstlisting}

\begin{tabular}{@{}lcccc@{}}
\toprule
\textbf{Profile} & $w_{\text{latency}}$ & $w_{\text{dollar}}$ & $w_{\text{energy}}$ & $w_{\text{queue}}$ \\
\midrule
Balanced & 1.0 & 1.0 & 1.0 & 1.0 \\
EnergyConstrained & 0.5 & 0.5 & 3.0 & 1.0 \\
CostSensitive & 0.5 & 3.0 & 0.5 & 1.0 \\
LatencyCritical & 3.0 & 0.1 & 0.1 & 2.0 \\
\bottomrule
\end{tabular}

Profiles are set at fleet level. Individual interests can override the fleet profile:

\begin{lstlisting}[language=Rust]
Interest {
    subject: "v2x.bsm.>",
    priority: Priority::High,
    cost_profile: Some(CostProfile::LatencyCritical),
    // ...
}
\end{lstlisting}

\subsection{Hybrid Logical Clocks}

Each device maintains a hybrid logical clock (HLC) consisting of:
\begin{itemize}
    \item $\text{wall} : \mathbb{N}$ (wall-clock time in milliseconds)
    \item $\text{counter} : \mathbb{N}$ (logical counter for sub-millisecond ordering)
    \item $\text{node} : D$ (node identifier for tie-breaking)
\end{itemize}

\textbf{HLC Operations:}

\begin{lstlisting}[language=Rust]
fn tick(hlc: &mut HLC) {
    let pt = physical_time();
    if pt > hlc.wall {
        hlc.wall = pt;
        hlc.counter = 0;
    } else {
        hlc.counter += 1;
    }
}

fn send(hlc: &mut HLC, message: &mut Message) {
    tick(hlc);
    message.hlc = hlc.clone();
}

fn receive(hlc_local: &mut HLC, message: &Message) {
    let hlc_remote = &message.hlc;
    let pt = physical_time();
    if pt > hlc_local.wall && pt > hlc_remote.wall {
        hlc_local.wall = pt;
        hlc_local.counter = 0;
    } else if hlc_remote.wall > hlc_local.wall {
        hlc_local.wall = hlc_remote.wall;
        hlc_local.counter = hlc_remote.counter + 1;
    } else if hlc_local.wall > hlc_remote.wall {
        hlc_local.counter += 1;
    } else {  // equal wall times
        hlc_local.counter = max(hlc_local.counter, hlc_remote.counter) + 1;
    }
}
\end{lstlisting}

\textbf{Total ordering:} Define $\text{hlc}_1 < \text{hlc}_2$ iff:
\begin{itemize}
    \item $\text{wall}(\text{hlc}_1) < \text{wall}(\text{hlc}_2)$, or
    \item $\text{wall}(\text{hlc}_1) = \text{wall}(\text{hlc}_2) \land \text{counter}(\text{hlc}_1) < \text{counter}(\text{hlc}_2)$, or
    \item $\text{wall}(\text{hlc}_1) = \text{wall}(\text{hlc}_2) \land \text{counter}(\text{hlc}_1) = \text{counter}(\text{hlc}_2) \land \text{node}(\text{hlc}_1) < \text{node}(\text{hlc}_2)$
\end{itemize}

\subsection{CRDT State Model}

Avena uses conflict-free replicated data types (CRDTs) to ensure convergence without coordination. Each state component uses a specific CRDT type:

\textbf{Device Registry:} OR-Map$\langle$DeviceId, DeviceInfo$\rangle$

Each DeviceInfo is a struct where fields use LWW-Register semantics with HLC timestamps:

\begin{lstlisting}[language=Rust]
struct DeviceInfo {
    id: DeviceId,
    public_key: Ed25519PublicKey,
    roles: LWWRegister<Set<Role>>,
    labels: LWWRegister<Map<String, String>>,
    position: LWWRegister<Option<GeoPosition>>,
    certificate: LWWRegister<Certificate>,
    // ... additional fields
}
\end{lstlisting}

\textbf{Merge operation:} For two DeviceInfo values with the same id:
\[
\text{merge}(a, b).\text{field} = 
\begin{cases}
a.\text{field} & \text{if } a.\text{field.hlc} > b.\text{field.hlc} \\
b.\text{field} & \text{otherwise}
\end{cases}
\]

\textbf{Workload Specifications:} OR-Map$\langle$WorkloadId, WorkloadSpec$\rangle$ with field-level LWW-Register merge.

\textbf{Active Interests:} OR-Set$\langle$Interest$\rangle$ with add-wins semantics:

\begin{lstlisting}[language=Rust]
struct InterestSet {
    // Each interest identified by (owner_device, interest_id)
    // Tombstones represent terminated interests
    entries: Map<(DeviceId, InterestId), InterestEntry>,
}

struct InterestEntry {
    interest: Interest,
    hlc: HLC,
    terminated: bool,  // true = tombstone
}
\end{lstlisting}

\textbf{Cost Metadata:} Map$\langle$DeviceId, LWWRegister$\langle$CostMetadata$\rangle\rangle$ (single-writer per device).

\textbf{Convergence Property:} Under $(T, i)$-connectivity, and assuming all clocks have bounded drift $\delta$, any gossip state update at time $\tau_0$ propagates to all devices by time $\tau_0 + T$. CRDT semantics ensure that all devices converge to the same state regardless of propagation order.

\subsection{Tombstone Management}

Terminated interests and deleted workloads are represented as tombstones rather than removed from the CRDT. This prevents ``resurrection'' where a delete fails to propagate to a replica that still holds the original entry.

\textbf{Tombstone lifecycle:}
\begin{enumerate}
    \item Entry marked as terminated with new HLC timestamp
    \item Tombstone propagates via gossip like any other update
    \item Tombstones retained for TTL period (default 6--12 months)
    \item After TTL, tombstones eligible for garbage collection
    \item GC events broadcast via gossip to coordinate cleanup
\end{enumerate}

\textbf{Resurrection edge case:} If a device rejoins the network after being offline longer than the tombstone TTL, it may attempt to re-propagate entries that were deleted. Defense: devices compare entry HLC against a ``minimum valid HLC'' threshold that advances with time. Entries with HLC older than the threshold are rejected.

\textbf{Formal GC condition:} An entry $e$ with tombstone timestamp $ts$ can be garbage collected at device $d$ at time $t$ if:
\[
t - ts > \text{TTL} \land \forall d' \in D : d' \text{ has received } e.\text{tombstone}
\]

The second condition is approximated by waiting for TTL to exceed expected maximum propagation time under worst-case connectivity.

\subsection{Local State Structure}

Each device $d$ maintains:

\begin{lstlisting}[language=Rust]
struct LocalState {
    // Gossip state (CRDT-based)
    gossip: GossipState,
    
    // Application message queues (not CRDT)
    messages: MessageState,
    
    // Local HLC
    clock: HLC,
}

struct GossipState {
    workloads: ORMap<WorkloadId, WorkloadSpec>,
    devices: ORMap<DeviceId, DeviceInfo>,
    cost_metadata: Map<DeviceId, LWWRegister<CostMetadata>>,
    interests: ORSet<Interest>,
    time_services: ORMap<DeviceId, TimeServiceInfo>,
}

struct MessageState {
    pending_outbound: PriorityQueue<(Subject, Message, InterestId)>,
    pending_inbound: PriorityQueue<(Message, InterestId)>,
    delivered: BoundedSet<(Subject, HLC)>,  // for deduplication
}
\end{lstlisting}

